%! Solving Rational Equations — Unit 4, Lesson 4.6 — Warm-Up
\documentclass[10pt]{article}
\usepackage{algebra2-article}
\usepackage{algebra2-boxes}
\newcommand{\TallMath}[1]{$\displaystyle #1\rule[-1.4em]{0pt}{3.2em}$}

\begin{document}

\pageheader{Unit 4, Lesson 4.6}{Warm-Up}
\namedateperiod

\vspace{0.10in}

\begin{notesbox}{Getting Ready: Skills We'll Use Today}
\small Three skills. The first two you already own. The third is a fact about \emph{equations} that
explains everything strange that happens today.

\vspace{4pt}
\textbf{1. Restrictions, off the original denominator} (Lesson 4.1). List every excluded value.
\begin{enumerate}[label=(\alph*), leftmargin=2.1em, itemsep=3pt]
  \item $\dfrac{5}{x-4}$: \; $x\neq\blank{1.4cm}$
  \item $\dfrac{x}{2x+6}$: \; $x\neq\blank{1.4cm}$
  \item $\dfrac{7}{x^2-9}$: \; $x\neq\blank{2.4cm}$
  \item $\dfrac{x+1}{x^2-2x-8}$: \; $x\neq\blank{2.4cm}$
\end{enumerate}

\vspace{4pt}
\textbf{2. The LCD, from factored denominators} (Lesson 4.3). Write the least common denominator.
\begin{enumerate}[label=(\alph*), leftmargin=2.1em, itemsep=3pt]
  \item $\dfrac{1}{x}$ and $\dfrac{3}{x-2}$: \quad LCD $=\blank{3.0cm}$
  \item $\dfrac{2}{x+3}$ and $\dfrac{5}{x^2-9}$: \quad LCD $=\blank{3.0cm}$
  \item $\dfrac{4}{x}$, \; $\dfrac{1}{x+1}$, \; and the number $3$: \quad LCD $=\blank{3.0cm}$
\end{enumerate}

\vspace{5pt}
\textbf{3. A fact about equations.} Start with a statement that is plainly \textbf{false}:
\[ 5 = 2 \]
\begin{enumerate}[label=(\alph*), leftmargin=2.1em, itemsep=4pt]
  \item Multiply both sides by $3$. New statement: \blank{2.4cm} \; Is it true or false?
        \blank{1.6cm} \\[2pt]
        Could you get back to $5=2$? Yes --- \blank{4.6cm} both sides by $3$.
  \item Now multiply both sides of $5=2$ by $\mathbf{0}$. New statement: \blank{2.4cm} \;
        True or false? \blank{1.6cm}
  \item Can you get back to $5=2$ from that? \blank{1.2cm} \; Why not? \blank{5.6cm}
  \item Finish the rule: multiplying both sides of an equation by a \emph{nonzero} number can always be
        undone, so it never changes which values make the equation true. But multiplying both sides by
        \blank{1.4cm} can turn a \blank{1.6cm} statement into a \blank{1.6cm} one.
\end{enumerate}
\end{notesbox}

\end{document}
